\chapter{Linear and Integer Programming}

A woodworking shop receives 24 sheets of plywood and 6 hours of assembly-machine time for the
day. One bookshelf uses 6 sheets of plywood and 1 hour of assembly. One bench uses 4 sheets
and 2 hours. A shelf earns a profit of 5 units, a bench 4 units. How many shelves and how many
benches should it make?

The first instinct picks shelves, because the profit per unit is higher. That instinct is
wrong: one shelf consumes as much plywood as one and a half benches. The best answer of the
continuous model is 3 shelves and 1.5 benches, with profit 21, and half a bench cannot be
sold. Rounding down gives 19; 4 shelves without benches gives 20. So the best integer answer
is not the rounding of the best fractional answer, and the difference is an error of reasoning,
not an error of arithmetic.

The tools of the previous chapters are not enough. Chapter 8 taught us to solve
$A\mathbf{x}=\mathbf{b}$, but here there are fewer equations than variables, so there are
infinitely many answers and we must choose the best one. Chapter 23 gives the KKT conditions,
but in this problem the gradient of the objective is constant and never vanishes: the answer
clings to the boundary of the feasible region, not to a stationary point. The guarantee of
Chapter 21 that a local optimum is also global does not show where the optimum is, and it
disappears as soon as the variables are required to be integers.

For that reason this chapter takes three steps: writing the problem as a linear program and
recognizing that its answer lies at a corner of the feasible region; putting a price on every
constraint through duality, so that optimality can be checked without sweeping the whole
region; then facing the requirement of integrality, when the continuous model gives a valid
bound that is not necessarily attained.

\begin{untukanak}
Imagine you have 24 sheets of colored paper and 6 hours to make two kinds of classroom
decoration. The first decoration uses 6 sheets and 1 hour, the second uses 4 sheets and 2
hours. The first decoration is in greater demand, so its price is higher. How many of each
should you make?

If you try things out, there are many mixtures that fit. What is surprising is that the most
profitable mixture turns out to use half a decoration. Half a decoration cannot be sold, and
this is where common sense begins to deceive: if you round down, the result is not the best
one. There is another mixture, entirely whole, that is in fact more profitable, even though it
does not look like the rounding of anything.

The question that makes one curious is this. If there are thousands of possible mixtures, how
do we know that one mixture is the best without trying them all? And if your friend offers one
extra sheet of paper, what is a fair price for that sheet? In this chapter we will learn why
the best answer always lies at a corner of the possible region, why every material has a price
that can be computed, and why a fractional answer must be rounded with care, not with instinct.
\end{untukanak}

\section{From a production plan to standard form}
\tingkat{basic}

A production plan is a pair of numbers: how many shelves and how many benches. A plan is called
feasible when its use of materials does not exceed the supply. For $x_{1}$ shelves and
$x_{2}$ benches, the conditions are
\[
  6x_{1}+4x_{2}\le24,\qquad x_{1}+2x_{2}\le6,\qquad x_{1}\ge0,\quad x_{2}\ge0 .
\]
Each condition cuts the plane with one line and leaves one side; the intersection is a
quadrilateral with four corners. The profit is $5x_{1}+4x_{2}$, and we want that number as
large as possible.

The pattern of this kind of problem is called a linear program. The name is apt: everything
involved is a straight line, with no curves and no powers.

\kamus{linear program}{program linear}{A problem of finding the best value of a quantity
computed by a weighted sum, when all the rules take the shape of straight lines.}

\begin{definisi}[Linear program]
A linear program is the problem of maximizing or minimizing a linear function
$\mathbf{c}^{\top}\mathbf{x}$ over variables $\mathbf{x}\in\mathbb{R}^{n}$ subject to
constraints in the form of linear inequalities or equalities. The set of points satisfying all
the constraints is called the \emph{feasible region}. If the feasible region is empty, the
program is called \emph{infeasible}; if the objective value can be made as large as we please,
the program is called \emph{unbounded}.\index{linear program}
\end{definisi}

Software and textbooks use one uniform written pattern, and that is the pattern used
throughout this chapter.

\begin{definisi}[Standard form]
A linear program in \emph{standard form} is
\[
  \text{maximize } \mathbf{c}^{\top}\mathbf{x}
  \quad\text{with}\quad A\mathbf{x}=\mathbf{b},\qquad \mathbf{x}\ge\mathbf{0},
\]
where $A$ is $m\times n$ of rank $m$, $\mathbf{b}\in\mathbb{R}^{m}$, and
$\mathbf{c}\in\mathbb{R}^{n}$. Every linear program can be brought to this form without
changing its optimal value.\index{standard form}
\end{definisi}

Four translations do all that work. A constraint
$\mathbf{a}^{\top}\mathbf{x}\le b$ becomes an equality by adding a variable that records the
unused remainder; a constraint $\mathbf{a}^{\top}\mathbf{x}\ge b$ by subtracting one. A
variable that may take negative values is replaced by the difference of two nonnegative
variables, and maximizing $\mathbf{c}^{\top}\mathbf{x}$ is the same as minimizing
$-\mathbf{c}^{\top}\mathbf{x}$.

\kamus{slack variable}{variabel kendur}{A number recorded as the unused remainder, so that
``no more than'' turns into ``equal to''.}

For the shop above, the standard form uses $s_{1}$ as the plywood remainder and $s_{2}$ as the
remainder of assembly hours:
\[
  \begin{array}{rcl}
    6x_{1}+4x_{2}+s_{1} &=& 24,\\
    x_{1}+2x_{2}+s_{2} &=& 6,
  \end{array}
  \qquad x_{1},x_{2},s_{1},s_{2}\ge0 .
\]
Here $m=2$ and $n=4$, so $A$ is $2\times4$. The two slack variables are not decoration: later
both appear as columns in the simplex tableau and as a sign of which constraint is still loose.

Standard form also demands $\mathbf{b}\ge\mathbf{0}$, and that can always be arranged by
multiplying one row by $-1$. What cannot be arranged is the rank: rows of $A$ that are
linearly dependent must be discarded first, since otherwise the basis matrix chosen may be
singular.

\section{Geometry: why the answer is at a corner}
\tingkat{intermediate}

The four corners of that quadrilateral are $(0;0)$, $(4;0)$, $(3;1.5)$, and $(0;3)$, with
profits $0$, $20$, $21$, and $12$. The best is at a corner, and that is no coincidence: in a
linear program the best value can always be found at a corner, which is why the simplex method
only travels from corner to corner.

\begin{figure}[!ht]
\centering
\begin{tikzpicture}
\begin{axis}[width=9.3cm,height=6.3cm,
  xlabel={$x_{1}$ (shelves)},ylabel={$x_{2}$ (benches)},
  xmin=-0.35,xmax=5.1,ymin=-0.35,ymax=3.9,
  axis lines=left,xtick={0,1,2,3,4,5},ytick={0,1,2,3},
  tick label style={font=\scriptsize},label style={font=\scriptsize}]
\addplot[forget plot,fill=layerteal!10,draw=layerteal,thick]
  coordinates {(0,0) (4,0) (3,1.5) (0,3)} \closedcycle;
\addplot[forget plot,domain=0:2.4,densely dotted,thick,layerindigo] {3-1.25*x};
\addplot[forget plot,domain=0:3.6,densely dotted,thick,layerindigo] {4.5-1.25*x};
\addplot[forget plot,domain=0.2:4.2,thick,accent] {5.25-1.25*x};
\addplot[forget plot,only marks,mark=*,mark size=1.7pt,draw=ink]
  coordinates {(0,0) (4,0) (0,3)};
\addplot[forget plot,only marks,mark=*,mark size=2.6pt,draw=accent]
  coordinates {(3,1.5)};
\node[font=\scriptsize,anchor=north west] at (axis cs:0,0) {$(0;0)$};
\node[font=\scriptsize,anchor=north east] at (axis cs:4,0) {$(4;0)$};
\node[font=\scriptsize,anchor=south west] at (axis cs:0,3) {$(0;3)$};
\node[font=\scriptsize,anchor=west] at (axis cs:3.05,1.45) {$(3;1.5)$};
\node[font=\scriptsize,anchor=west,layerindigo] at (axis cs:0.15,1.15) {$z=12$};
\node[font=\scriptsize,anchor=west,layerindigo] at (axis cs:0.15,2.15) {$z=18$};
\node[font=\scriptsize,anchor=west,accent] at (axis cs:0.15,3.45) {$z=21$};
\draw[-{Stealth},thick,accent] (axis cs:2.3,2.75) -- (axis cs:3.25,3.35)
  node[font=\scriptsize,above,accent,pos=0.9] {profit rises};
\end{axis}
\end{tikzpicture}
\caption{The feasible region of the shop and three lines of equal profit. The polygon has four
vertices, $(0;0)$, $(4;0)$, $(3;1.5)$, $(0;3)$, with values $0$, $20$, $21$, $12$. The lines
$z=12$ and $z=18$ still cut the feasible region; the line $z=21$ only touches the vertex
$(3;1.5)$, so no feasible point has a higher value.}
\label{fig:ch22-daerah-optimum}
\end{figure}

Figure~\ref{fig:ch22-daerah-optimum} sharpens that coincidence. The three lines there are
lines of equal profit, $5x_{1}+4x_{2}=12$, $18$, and $21$. A higher profit line is farther from
the origin, and shifting it in that direction moves it out of the feasible region; the line
$21$ only touches the vertex $(3;1.5)$.

A corner has a more precise name, and the name does not refer to a picture.

\begin{definisi}[Vertex]
Let $P\subseteq\mathbb{R}^{n}$ be a convex set. A point $\mathbf{x}\in P$ is called a
\emph{vertex} or an \emph{extreme point} if there are no $\mathbf{y},\mathbf{z}\in P$ with
$\mathbf{y}\ne\mathbf{z}$ and $t\in(0,1)$ such that
$\mathbf{x}=t\mathbf{y}+(1-t)\mathbf{z}$. In other words, $\mathbf{x}$ is not the midpoint of
a segment between two other points of $P$.\index{vertex}
\end{definisi}

In the quadrilateral above, the point $(3;1.5)$ is a vertex: every segment inside the
quadrilateral that passes through it must end at that point itself. Conversely $(2;1)$ is not a
vertex, since it is the midpoint between $(1;1)$ and $(3;1)$.

\kamus{vertex}{verteks}{A corner of the feasible region, namely a point that never becomes the
midpoint between two other feasible points.}

The first step of the proof turns the negation into a direction that can be used.

\begin{lemma}[A point that is not a vertex gives a line]
If $\mathbf{x}\in P$ is not a vertex, then there is $\mathbf{d}\ne\mathbf{0}$ such that
$\mathbf{x}+\mathbf{d}\in P$ and $\mathbf{x}-\mathbf{d}\in P$.
\end{lemma}

\begin{proof}
Since $\mathbf{x}$ is not a vertex, there are $\mathbf{y}\ne\mathbf{z}\in P$ and $t\in(0,1)$
with $\mathbf{x}=t\mathbf{y}+(1-t)\mathbf{z}$. Set
$s=\tfrac12\min\{t,1-t\}>0$ and $\mathbf{d}=s(\mathbf{y}-\mathbf{z})$. Then
$\mathbf{x}+\mathbf{d}=(t+s)\mathbf{y}+(1-t-s)\mathbf{z}$ and
$\mathbf{x}-\mathbf{d}=(t-s)\mathbf{y}+(1-t+s)\mathbf{z}$. Since $0<t\pm s<1$, both
points are convex combinations of $\mathbf{y}$ and $\mathbf{z}$, hence members of $P$ because
$P$ is convex. Also $\mathbf{d}\ne\mathbf{0}$, since $\mathbf{y}\ne\mathbf{z}$ and $s>0$.
\end{proof}

Now the statement that makes the simplex method reasonable. In everyday language: if the
feasible region is enclosed and the objective does not grow in any direction, then the highest
summit is at one of the corners, not in the middle of the plane.

\begin{teorema}[The optimum is attained at a vertex]
Let $P=\{\mathbf{x}:A\mathbf{x}\le\mathbf{b}\}$ be nonempty and bounded, and suppose
$\mathbf{c}^{\top}\mathbf{x}$ is bounded above on $P$. Then there is a vertex
$\mathbf{x}^{*}$ of $P$ with $\mathbf{c}^{\top}\mathbf{x}^{*}
=\max_{\mathbf{x}\in P}\mathbf{c}^{\top}\mathbf{x}$.
\index{optimum at a vertex}
\end{teorema}

\begin{proof}
The optimal set $F=\{\mathbf{x}\in P:\mathbf{c}^{\top}\mathbf{x}=p^{*}\}$ is nonempty, closed,
bounded, and convex, so $F$ is also a polytope. Write the constraints of $P$ as
$\mathbf{a}_{i}^{\top}\mathbf{x}\le b_{i}$, $i=1,\dots,k$. Choose $\mathbf{x}^{*}\in F$
that maximizes the number of active constraints, namely the constraints with
$\mathbf{a}_{i}^{\top}\mathbf{x}^{*}=b_{i}$.

Suppose $\mathbf{x}^{*}$ is not a vertex. By the lemma there is $\mathbf{d}\ne\mathbf{0}$
with $\mathbf{x}^{*}\pm\mathbf{d}\in F$. Every constraint active at
$\mathbf{x}^{*}$ is also active at both of those points, so
$\mathbf{a}_{i}^{\top}\mathbf{d}=0$ for all active $i$. Consider the line
$\mathbf{x}^{*}+s\mathbf{d}$ for $s\in\mathbb{R}$. Its intersection with $P$ is a closed
interval $[s_{-},s_{+}]$ that contains $[-1,1]$, and because $P$ is bounded both endpoints are
finite. The endpoint $s_{+}$ is determined by the constraint with $\mathbf{a}_{i}^{\top}\mathbf{d}>0$
that is reached first; that constraint is active at $s_{+}$ but not at $\mathbf{x}^{*}$, so the
new point is feasible and active at more constraints. Its value stays $p^{*}$ because
$\mathbf{c}^{\top}(\mathbf{x}^{*}+s_{+}\mathbf{d})=p^{*}+s_{+}\mathbf{c}^{\top}
\mathbf{d}=p^{*}$, and $\mathbf{c}^{\top}\mathbf{d}=0$ since both points
$\mathbf{x}^{*}\pm\mathbf{d}$ have value $p^{*}$. That new point is a member of $F$ with more
active constraints, contradicting the choice of $\mathbf{x}^{*}$; hence $\mathbf{x}^{*}$ is a
vertex.
\end{proof}

\begin{bacakeras}
Two voices, read alternately with a friend. Tap the table once each time the word ``summit'' is
spoken.

``Where is the best answer?'' ``At the summit.''
``Why not in the middle?'' ``Because in the middle there is a neighbor on two sides.''
``If there is a neighbor, why is that wrong?'' ``Because there is still a direction to move.''
``Move in which direction?'' ``In the direction that climbs.''
``If there is no climbing direction?'' ``Then we are already standing at the summit.''

Repeat the three middle lines, slower and slower, with the same taps: in the middle there is a
neighbor, in the middle one can still move, a summit has no climbing direction.

The line the reader must repeat, once, in a clear voice: the best answer is at the summit, not
in the middle, and a summit is recognized by the absence of a climbing direction.
\end{bacakeras}

That proof uses the boundedness of $P$. If $P$ is unbounded but the objective is bounded above,
the conclusion is still true; the route goes through the decomposition of $P$ into the convex
hull of its vertices plus the cone of recession directions, and the complete proof is in
\texttt{WSPrinted.pdf}, in the section on the structure of polyhedral sets.

The number of vertices is finite, but not always small. In standard form, a vertex corresponds
to a choice of $m$ linearly independent columns of $A$: the chosen columns determine
$\mathbf{x}_{B}=B^{-1}\mathbf{b}$, the other columns are zero, and the point is feasible when
$\mathbf{x}_{B}\ge\mathbf{0}$. There are $\binom{n}{m}$ candidates, which for $n=100$ and
$m=50$ exceeds $10^{29}$; checking them all is not a sensible plan. This is where the simplex
method comes in: it walks from a vertex to a better neighboring vertex.

\begin{tangan}
Materials: a tiled floor, two lengths of string or paper tape, and a handful of buttons.

Stretch two straight strings as two boundaries; one for plywood, one for assembly hours. Stand
on a corner tile and call that tile the origin: each tile to the right adds one shelf, each
tile forward adds one bench.

Walk along the strings and stop only at tiles where two strings meet. Put a button there, and
only there. What must be visible: all the buttons cling to the turns of the fence, not one of
them stands in the middle of the yard. After that, count the profit at each button on your
fingers, then move one button to the most profitable corner.

Shift one string one tile inward, then repeat: the buttons move, the winner is still the
button at the corner.
\end{tangan}

\section{The simplex method: one tableau, two pivots}
\tingkat{intermediate}

The idea is simple. Stand at one vertex, check whether any edge direction raises the profit; if
one does, walk to the vertex at its end. If none does, stop: no neighbor is better, which in a
linear program means no other point is better.

That edge direction is determined by the \emph{reduced cost}: how much the profit increases
when a variable that is now zero is raised by one unit, at the optimal expense of the basis
variables.

\kamus{basic feasible solution}{solusi layak basis}{A point obtained by setting some variables
to zero and then solving for the rest, which happens to still satisfy all the bounds; such a
point is always a vertex.}

For the shop above, start from the vertex $(0;0)$, namely $s_{1}=24$ and $s_{2}=6$. Its first
tableau contains the profit row $z$ and one row for each constraint:
\begin{center}
\begin{tabular}{c|cccc|c}
\hline
basis & $x_{1}$ & $x_{2}$ & $s_{1}$ & $s_{2}$ & right-hand side \\
\hline
$z$   & $-5$ & $-4$ & $0$ & $0$ & $0$ \\
$s_{1}$ & $6$ & $4$ & $1$ & $0$ & $24$ \\
$s_{2}$ & $1$ & $2$ & $0$ & $1$ & $6$ \\
\hline
\end{tabular}
\end{center}
The $z$ row is written as $z-5x_{1}-4x_{2}=0$; a negative coefficient means that raising that
variable raises $z$, so the column with the most negative coefficient enters, here $x_{1}$ with
$-5$.

The ratio test decides who leaves: divide the right-hand side by the positive entries in the
entering column, then take the smallest ratio. $24/6=4$ and $6/1=6$, so $s_{1}$ leaves and the
pivot is $6$. Dividing the pivot row by $6$ and then clearing that column from the other rows
gives
\[
  \begin{array}{rcl}
    x_{1}+\tfrac{2}{3}x_{2}+\tfrac{1}{6}s_{1} &=& 4,\\
    \tfrac{4}{3}x_{2}-\tfrac{1}{6}s_{1}+s_{2} &=& 2,\\
    z-\tfrac{2}{3}x_{2}+\tfrac{5}{6}s_{1} &=& 20 .
  \end{array}
\]
The point is now $(4;0)$ with profit $20$, and a negative coefficient remains: $x_{2}$ with
$-2/3$. The ratio test gives $4/(2/3)=6$ and $2/(4/3)=1.5$, so $s_{2}$ leaves and the pivot is
$4/3$. After one more round,
\[
  \begin{array}{rcl}
    x_{1}+\tfrac{1}{4}s_{1}-\tfrac{1}{2}s_{2} &=& 3,\\
    x_{2}-\tfrac{1}{8}s_{1}+\tfrac{3}{4}s_{2} &=& \tfrac{3}{2},\\
    z+\tfrac{3}{4}s_{1}+\tfrac{1}{2}s_{2} &=& 21 .
  \end{array}
\]
The $z$ row now has no negative coefficient, so no direction raises the profit:
$(3;1.5)$ is optimal with profit $21$. All the arithmetic above was rechecked with exact
fractions, so that no decimal rounding slipped in.

That last row is more than a stopping sign. The coefficients of $s_{1}$ and
$s_{2}$ there, $3/4$ and $1/2$, are the prices of the two materials. Those numbers come from the
basis matrix, $\mathbf{y}^{\top}=\mathbf{c}_{B}^{\top}B^{-1}$, and
$\mathbf{b}^{\top}\mathbf{y}=24\cdot\tfrac34+6\cdot\tfrac12=21$, the same as the profit.
That equality is duality, and the next part explains why it is itself already
a proof of optimality.

Two remarks. The most-negative-coefficient rule does not guarantee termination, since there are
small examples with an endless cycle; Bland's rule, which picks the smallest index, guarantees
termination. Modern software uses the dual simplex or interior-point methods, but reports the
same vertices and the same prices.

\section{Duality: a price for every constraint}
\tingkat{advanced}

Every scarce material has a hidden price. One extra sheet of plywood raises the profit, and the
size of that rise is the price of the sheet to the shop. That price is not written in the
original problem, but it can be computed, and duality is the way to compute it.

Multiply the plywood constraint by $y_{1}\ge0$ and the hours constraint by $y_{2}\ge0$, then
add them. The result is
\[
  (6y_{1}+y_{2})x_{1}+(4y_{1}+2y_{2})x_{2}\le24y_{1}+6y_{2}.
\]
If $6y_{1}+y_{2}\ge5$ and $4y_{1}+2y_{2}\ge4$, the left-hand side is at least
$5x_{1}+4x_{2}$, which is the profit. So $24y_{1}+6y_{2}$ becomes an upper bound on the profit
for every price satisfying those two conditions, and the best bound is made as small as
possible.

\begin{definisi}[Primal and dual pair]
For the linear program
\[
  \text{(P)}\quad \text{maximize } \mathbf{c}^{\top}\mathbf{x}
  \ \text{ with } A\mathbf{x}\le\mathbf{b},\ \mathbf{x}\ge\mathbf{0},
\]
\emph{its dual program} is
\[
  \text{(D)}\quad \text{minimize } \mathbf{b}^{\top}\mathbf{y}
  \ \text{ with } A^{\top}\mathbf{y}\ge\mathbf{c},\ \mathbf{y}\ge\mathbf{0}.
\]
The variables $\mathbf{y}$ are called \emph{dual variables} or \emph{shadow prices}.
\index{duality}\index{shadow price}
\end{definisi}

For the shop above, the dual minimizes $24y_{1}+6y_{2}$ subject to
$6y_{1}+y_{2}\ge5$, $4y_{1}+2y_{2}\ge4$, and $y_{1},y_{2}\ge0$. Its solution is
$y_{1}=3/4$, $y_{2}=1/2$, with value $21$, the same as the best profit.

That the upper bound from the dual never falls below the primal value is the first thing to
hold on to, and its proof is three lines.

\begin{teorema}[Weak duality]
For every feasible $\mathbf{x}$ of (P) and every feasible $\mathbf{y}$ of (D) we have
$\mathbf{c}^{\top}\mathbf{x}\le\mathbf{b}^{\top}\mathbf{y}$.
\index{weak duality}
\end{teorema}

\begin{proof}
Since $\mathbf{y}\ge\mathbf{0}$ and $\mathbf{b}-A\mathbf{x}\ge\mathbf{0}$, we have
$\mathbf{b}^{\top}\mathbf{y}\ge\mathbf{y}^{\top}A\mathbf{x}$. Since
$\mathbf{x}\ge\mathbf{0}$ and $A^{\top}\mathbf{y}-\mathbf{c}\ge\mathbf{0}$, we have
$\mathbf{y}^{\top}A\mathbf{x}\ge\mathbf{c}^{\top}\mathbf{x}$. The two chains together give the
required result.
\end{proof}

The practical consequence is large: a pair of feasible points, one primal and one dual, with
the same value already proves both optimal without inspecting any other point. In the shop
above, $(3;1.5)$ and $(0.75;\,0.5)$ give the value $21$, so $21$ can be neither improved nor
lowered.

What weak duality does not give is a guarantee that the two values always meet. That guarantee
is the content of the next theorem.

\begin{teorema}[Strong duality]
If (P) or (D) has a finite optimal value, then both have optimal values and those values are
equal.
\end{teorema}

The proof is not carried out here. There are two routes: through Farkas' lemma on systems of
linear inequalities, or through the fact that the simplex stops at a tableau with no negative
coefficient, and that tableau at the same time gives a feasible dual $\mathbf{y}$ of the same
value. The complete proof is in \texttt{WSPrinted.pdf}, in the section on linear programming
duality, and in Boyd and Vandenberghe, section 5.2.

From that equality comes the condition used to check a computation quickly. In everyday
language: a constraint that is still loose does not need to be paid for, and a variable that is
used must pay its price exactly.

\begin{akibat}[Complementary slackness]
If $\mathbf{x}$ and $\mathbf{y}$ are both optimal, then
\[
  y_{i}\bigl(b_{i}-\mathbf{a}_{i}^{\top}\mathbf{x}\bigr)=0 \ \text{ for every } i,
  \qquad
  x_{j}\bigl((A^{\top}\mathbf{y})_{j}-c_{j}\bigr)=0 \ \text{ for every } j .
\]
Conversely, if $\mathbf{x}$ is feasible for (P), $\mathbf{y}$ is feasible for (D), and both
groups of equalities hold, then both are optimal.
\index{complementary slackness}
\end{akibat}

\begin{proof}
Since both are optimal, $\mathbf{c}^{\top}\mathbf{x}=\mathbf{b}^{\top}\mathbf{y}$, and the
chain in the proof of weak duality becomes a chain of equalities. Every term summed inside it is
nonnegative, so every term is zero; those terms are exactly
$y_{i}(b_{i}-\mathbf{a}_{i}^{\top}\mathbf{x})$ and
$x_{j}((A^{\top}\mathbf{y})_{j}-c_{j})$. The converse direction: if all terms are zero, the same
chain gives equality of values, and by weak duality that value is optimal for both problems.
\end{proof}

Now the shadow prices themselves. If the right-hand side is perturbed, the optimal value
changes, and the rate of change is the component of $\mathbf{y}$. Adding one sheet of plywood
raises the profit from $21$ to $21.75$; adding one assembly hour raises it to
$21.5$. So a sheet of plywood is worth $0.75$ units of profit, an assembly hour
$0.5$.

That price does not hold forever, and its range of validity can be computed. As long as the set
of active constraints does not change, $\mathbf{x}_{B}=B^{-1}\mathbf{b}$ is linear in
$\mathbf{b}$, so its value is linear too. For the shop above,
\[
  B^{-1}=\begin{pmatrix}\tfrac14&-\tfrac12\\[2pt]-\tfrac18&\tfrac34\end{pmatrix},
  \qquad
  \mathbf{x}_{B}=\begin{pmatrix}\tfrac14b_{1}-\tfrac12b_{2}\\[2pt]
  -\tfrac18b_{1}+\tfrac34b_{2}\end{pmatrix}.
\]
With $b_{2}=6$ fixed, the condition $\mathbf{x}_{B}\ge\mathbf{0}$ gives $b_{1}\in[12,36]$;
outside that interval one basis variable becomes negative and the slope changes. With $b_{1}=24$
fixed, the interval is $b_{2}\in[4,12]$; a shadow price is therefore always
reported together with its interval, not as a bare number.

One warning is useful from the start: shadow prices are an idea of the continuous model. In an
integer model, one extra sheet of plywood may raise the profit by $1$ or not at all, and there is
no single price that summarizes both.

\section{The LP relaxation: a valid bound, a deceptive rounding}
\tingkat{intermediate}

If shelves and benches may be fractional, the problem above is finished: $3$ shelves and
$1.5$ benches. In reality they may not, and the relation between the continuous problem and
the integer problem must be stated with care.

\begin{definisi}[Integer program, relaxation, and integrality gap]
An \emph{integer program} is a linear program in which some or all of the variables are
required to be integer. If all the variables are required to be integer, it is called a pure
integer program. Its \emph{LP relaxation} is the same linear program without the integrality
requirement. For a maximization problem with optimal values $z^{*}_{\mathrm{LP}}$ and
$z^{*}_{\mathrm{IP}}$, the \emph{integrality gap} is $z^{*}_{\mathrm{LP}}-z^{*}_{\mathrm{IP}}$.
\index{integer program}\index{relaxation}\index{integrality gap}
\end{definisi}

\kamus{integrality gap}{celah integralitas}{The distance between the bound that the continuous
model promises and the actual integer answer; that distance measures how far the continuous
model misleads.}

The reason the relaxation is useful is one sentence: the set of feasible integer points is part
of the set of feasible continuous points, so maximizing over the larger set cannot give a
smaller value.

\begin{proposisi}[The relaxation gives a valid bound]
For a maximization problem, $z^{*}_{\mathrm{IP}}\le z^{*}_{\mathrm{LP}}$. For a minimization
problem, the inequality reverses.
\index{relaxation bound}
\end{proposisi}

\begin{proof}
Every integer point that is feasible for the integer program is also feasible for its
relaxation, since the relaxation only removes conditions, it does not add any. So the feasible
region of the integer program is contained in the feasible region of the relaxation. A maximum
over a larger set is not smaller than the maximum over its subset.
\end{proof}

The bound is valid, but not always tight, and the difference can be very misleading when it is
used for rounding. Rounding fails in two ways, and both appear in our small example.

The first way: the rounded result is infeasible. The point $(3;1.5)$ rounded up becomes
$(3;2)$. Its profit is $23$, higher than $21$, but its plywood is
$6\cdot3+4\cdot2=26$ sheets, exceeding $24$: that is not a production plan, only a row of
numbers.

The second way, subtler: the rounded result is feasible but not the best. Rounding down gives
$(3;1)$ with profit $19$, although $(4;0)$ has profit $20$ and is entirely feasible.
Rounding the coordinates one by one in the nearest direction therefore gives a worse answer
than another integer point that is farther from $(3;1.5)$.

Both failures are a consequence of the shape of the feasible region. The set of feasible
integer points is not convex: $(3;1)$ and $(4;0)$ are members, but their midpoint
$(3.5;\,0.5)$ is not, because it is not integer. On a nonconvex set the vertex-optimum theorem
does not apply: no vertex is guaranteed to contain the answer, and no neighbor can be relied
upon.

\begin{figure}[!ht]
\centering
\begin{tikzpicture}
\begin{axis}[width=9.3cm,height=6.3cm,
  xlabel={$x_{1}$ (shelves)},ylabel={$x_{2}$ (benches)},
  xmin=-0.35,xmax=5.1,ymin=-0.35,ymax=3.9,
  axis lines=left,xtick={0,1,2,3,4,5},ytick={0,1,2,3},
  tick label style={font=\scriptsize},label style={font=\scriptsize}]
\addplot[forget plot,fill=layerteal!8,draw=layerteal,thick]
  coordinates {(0,0) (4,0) (3,1.5) (0,3)} \closedcycle;
\addplot[forget plot,only marks,mark=*,mark size=1.5pt,draw=layerindigo]
  coordinates {(0,0) (1,0) (2,0) (3,0) (4,0)
               (0,1) (1,1) (2,1) (3,1)
               (0,2) (1,2) (2,2) (0,3)};
\addplot[forget plot,only marks,mark=square*,mark size=2.4pt,draw=layerjade]
  coordinates {(4,0)};
\addplot[forget plot,only marks,mark=*,mark size=1.5pt,draw=layerplum]
  coordinates {(3,1)};
\addplot[forget plot,only marks,mark=o,mark size=3pt,thick,accent]
  coordinates {(3,1.5)};
\draw[densely dashed,layergrey] (axis cs:3,1) -- (axis cs:3,2.35);
\node[font=\scriptsize,anchor=west,layerjade] at (axis cs:4.1,-0.05) {$(4;0)$, $z=20$};
\node[font=\scriptsize,anchor=east,layerplum] at (axis cs:2.9,0.85) {$(3;1)$, $z=19$};
\node[font=\scriptsize,anchor=west,accent] at (axis cs:3.15,1.6) {$(3;1.5)$, $z=21$};
\node[font=\scriptsize,anchor=north,layergrey] at (axis cs:3,-0.05) {};
\end{axis}
\end{tikzpicture}
\caption{The integer lattice versus the relaxation point. There are $13$ feasible integer points,
marked with filled dots. The relaxation gives $(3;1.5)$ with value $21$, which stands midway
between $(3;1)$ with value $19$ and $(3;2)$, which is infeasible since
$6\cdot3+4\cdot2=26>24$. The best integer answer $(4;0)$ with value $20$ is in fact farther
from the relaxation point.}
\label{fig:ch22-kisi-relaksasi}
\end{figure}

Figure~\ref{fig:ch22-kisi-relaksasi} contains everything: thirteen feasible integer points, and
one relaxation point standing between two points. The one in between is not an integer point,
and the nearest is not necessarily the best.

Rounding also has no single meaning in many dimensions. With two variables, ``nearest integer
point'' can still be pictured; with a hundred variables, the nearest is not necessarily
feasible, and the feasible one is not necessarily the best. The practice is not to round, but to
use the relaxation bound to prune the search.

\begin{tantangan}
Maximize $6x_{1}+5x_{2}$ subject to $2x_{1}+3x_{2}\le12$, $4x_{1}+x_{2}\le13$, and
$x_{1},x_{2}\ge0$ integer. The relaxation gives $(2.7;\,2.2)$ with value $27.2$.
Rounding to the nearest integer, $(3;2)$, gives value $28$: higher than the value of any
relaxation point that is still feasible. Show why $28$ cannot be trusted, then determine the
correct integer answer together with its value.

Hint 1: check feasibility, not only the value. A higher number may come from a point that is
not a plan.

Hint 2: squeeze $x_{1}$. The largest integer that is still feasible is $3$; for
$x_{1}=3$, compute the bound on $x_{2}$ from both constraints and take the smaller one.

If you get stuck, go back to the section The LP relaxation: a valid bound, a deceptive
rounding.
\end{tantangan}

\section{Why integer programs are harder}
\tingkat{intermediate}

Up to here the question is fair: if the relaxation only adds one requirement, why is an integer
program not solved by solving its relaxation and then rounding? The short answer: the difficulty
is not in finding an answer, but in making sure that the answer is the best one.

Three differences explain it.

First, there is no simple certificate of optimality. In a linear program, a pair $\mathbf{x}$
and $\mathbf{y}$ of equal value is a finished proof that can be checked in time proportional to
the size of the data. In an integer program, the prices from the relaxation give only a bound,
not an integer point that attains it: the dual of the relaxation gives a lower bound, not an
answer.

Second, the difference between bound and answer can be large. In the vertex cover of the
complete graph on $n$ vertices, the relaxation allows $x_{v}=1/2$ for all $v$ with value $n/2$,
whereas the integer answer is $n-1$; the ratio $2(n-1)/n$ approaches $2$. The complete
computation is in the computational lab and in Exercise 10.

Third, the computational cost is of a different class, and the number of candidates grows
exponentially. A linear program is solved in time polynomial in the input size; Khachiyan's
result through the ellipsoid method and Karmarkar's through interior-point methods are the two
landmarks. The decision version of the integer program, ``is there a feasible integer point of
value at least $K$'', is NP-complete, and Karp's list of twenty-one problems contains knapsack,
which by itself has $2^{n}$ subsets. No polynomial algorithm is known for integer programming,
and if one existed, many other things would collapse with it.

\begin{untukanak}
Have you ever laid tiles on a floor? You measure the length of the room, then divide by the size
of the tile. The result is rarely a whole number. If you round down, there is a gap at the end
of the floor. If you round up, you cut a tile and throw away the rest. And sometimes the best
answer is neither of the two, but a laying pattern that is altogether different: start from the
middle, or use cut tiles on both sides.

Someone who looks only at the single number that comes out of the division will not find that
pattern. He needs two bounds: at most how many tiles are used, and at least how many. With those
two numbers he stops thinking about patterns that are plainly too wasteful, and concentrates on
the ones that may still win.

In this chapter we will learn why the bound from the continuous model is always valid, why
rounding can give an answer that is feasible but not the best, and how a computer uses that
bound to prune millions of possibilities without inspecting them one by one.
\end{untukanak}

\section{Branch and bound}
\tingkat{advanced}

If the relaxation answer is fractional, we do not round it, but split the problem into two and
ask again with the relaxation on each piece. Every integer answer satisfies one of two mutually
exclusive possibilities, $x_{2}\le1$ or $x_{2}\ge2$, so solving both is the same as solving the
original problem.

The procedure has four steps. Solve the relaxation at one node. If the result is integer, record
it as a candidate answer and do not continue with that node. If it is fractional, choose one
fractional variable, $x_{j}=v$, then create two children: $x_{j}\le\lfloor v\rfloor$ and
$x_{j}\ge\lceil v\rceil$. Before inspecting a node, compare its relaxation bound with the best
candidate; if it is not better, discard that node together with all its descendants. Stop when
no node remains.

That discarding rule is valid, and the reason is one line.

\begin{proposisi}[Pruning discards no answer]
Let $z_{\mathrm{LP}}(N)$ be the relaxation value at node $N$ and $z_{\mathrm{best}}$ the best
candidate already recorded. If $z_{\mathrm{LP}}(N)\le z_{\mathrm{best}}$, then no feasible
integer point at $N$ is better than $z_{\mathrm{best}}$.
\end{proposisi}

\begin{proof}
Every integer point that is feasible at $N$ is also feasible for the relaxation of $N$, so its
value is at most $z_{\mathrm{LP}}(N)$ by the relaxation-bound proposition. If
$z_{\mathrm{LP}}(N)\le z_{\mathrm{best}}$, all those points have value at most
$z_{\mathrm{best}}$, so there is nothing left to inspect.
\end{proof}

The tree for the shop above is small, and can be written entirely as a table.
\begin{center}
\begin{tabular}{llcc}
\hline
node & added constraint & relaxation $\mathbf{x}$ & $z_{\mathrm{LP}}$ \\
\hline
root & -- & $(3;\,1.5)$ & $21$ \\
left & $x_{2}\le1$ & $(10/3;\,1)$ & $20.\overline{6}$ \\
right & $x_{2}\ge2$ & $(2;\,2)$ & $18$ \\
left-left & $x_{2}\le1,\ x_{1}\le3$ & $(3;\,1)$ & $19$ \\
left-right & $x_{2}\le1,\ x_{1}\ge4$ & $(4;\,0)$ & $20$ \\
\hline
\end{tabular}
\end{center}
The root gives $21$ with $x_{2}$ fractional. Branching on $x_{2}$ gives two children: left
$20.\overline{6}$ with $x_{1}=10/3$ still fractional, right $18$ and already integer. The right
node is recorded as a candidate. At the left node, branching on $x_{1}$ gives $(3;1)$ with value
$19$ and $(4;0)$ with value $20$; the second updates the candidate. No node remains, and the
answer is $20$.

\begin{figure}[!ht]
\centering
\begin{tikzpicture}[
  simpul/.style={draw=layergrey,fill=layergrey!6,rounded corners=1.5pt,
    inner sep=3pt,font=\scriptsize,align=center,text width=3.0cm},
  pangkas/.style={draw=accent,fill=accent!7,rounded corners=1.5pt,
    inner sep=3pt,font=\scriptsize,align=center,text width=3.0cm},
  panah/.style={-{Stealth},thick,draw=layergrey}]
\node[simpul] (akar) at (0,0) {root\\ $z_{\mathrm{LP}}=21$\\ $(3;\,1.5)$};
\node[simpul] (kiri) at (-3.6,-1.75) {$x_{2}\le1$\\ $z_{\mathrm{LP}}=20.\overline{6}$\\
  $(10/3;\,1)$};
\node[pangkas] (kanan) at (3.6,-1.75) {$x_{2}\ge2$\\ $z_{\mathrm{LP}}=18$\\ $(2;\,2)$ integer};
\node[pangkas] (kk) at (-5.1,-3.6) {$x_{1}\le3$\\ $z_{\mathrm{LP}}=19$\\ $(3;\,1)$ integer};
\node[simpul] (kn) at (-2.1,-3.6) {$x_{1}\ge4$\\ $z_{\mathrm{LP}}=20$\\ $(4;\,0)$ integer};
\draw[panah] (akar.south) -- (kiri.north);
\draw[panah] (akar.south) -- (kanan.north);
\draw[panah] (kiri.south) -- (kk.north);
\draw[panah] (kiri.south) -- (kn.north);
\node[font=\scriptsize,accent,anchor=north] at (3.6,-2.55) {pruned: $18<20$};
\node[font=\scriptsize,accent,anchor=north] at (-5.1,-4.45) {pruned: $19<20$};
\end{tikzpicture}
\caption{The branch-and-bound tree for the shop: five nodes. The root gives $z=21$
with $x_{2}$ fractional. The branch $x_{2}\ge2$ ends at $(2;\,2)$ with value $18$; the branch
$x_{2}\le1$ branches again into $(3;\,1)$ with value $19$ and $(4;\,0)$ with value $20$.
The two nodes labeled integer whose values are below $20$ have a red frame: their bounds are no
longer better than the best candidate, so both are pruned without further inspection.}
\label{fig:ch22-pohon-percabangan}
\end{figure}

Figure~\ref{fig:ch22-pohon-percabangan} draws the same tree. Two nodes there
have a red frame, namely $(2;\,2)$ with value $18$ and $(3;\,1)$ with value $19$: both are
integer, but their bounds are below the best candidate $20$, so both are discarded.

The right node with value $18$ would not need further inspection if it were
fractional: the bound $18$ is already below the candidate $20$, so its whole branch can be
discarded. On an example this small the saving is not noticeable; on a problem with hundreds of
variables, that rule is what separates finishing in a minute from not finishing at all, since
without good bounds branch and bound degenerates into $2^{n}$ nodes.

Two things must be said honestly. The worst-case time is still exponential, and there is no
guarantee that the tree is small. The choice of branching variable and the order of the nodes
also strongly affect the size of the tree; software uses heuristics, for example the variable
closest to one half. Both explain why an integer program can be far slower than a linear program
of the same size.

\section{Worked example}
\tingkat{intermediate}

Now all the numbers are carried to the end, so that every claim above can be
rechecked by hand or with Python.

\begin{resep}
Four steps solve a two-variable linear program graphically:
\begin{enumerate}
  \item Turn each constraint into a line by replacing the inequality sign with an equals sign,
        then draw the line through two axis intercepts.
  \item Determine the feasible side by testing one point, usually the origin; shade the wrong
        side so that the feasible region stays clean.
  \item Intersect every pair of lines, discard the intersection points that do not satisfy all
        the constraints, and what remains are the vertices.
  \item Compute the objective value at every vertex, then take the largest for maximization or
        the smallest for minimization.
\end{enumerate}
A warning about order: step 3 must not be carried out before step 2. Intersection points
computed earlier often lie outside the feasible region and their values are in fact higher than
the true answer; a nice number that appears before feasibility is checked is hard to throw away.
\end{resep}

\emph{First, the linear program of the shop.} Its feasible region is a quadrilateral with four
vertices, and all four are inspected:
\begin{center}
\begin{tabular}{lcc}
\hline
vertex & active constraints & profit $5x_{1}+4x_{2}$ \\
\hline
$(0;\,0)$ & $x_{1}=0,\ x_{2}=0$ & $0$ \\
$(4;\,0)$ & $6x_{1}+4x_{2}=24,\ x_{2}=0$ & $20$ \\
$(3;\,1.5)$ & $6x_{1}+4x_{2}=24,\ x_{1}+2x_{2}=6$ & $21$ \\
$(0;\,3)$ & $x_{1}+2x_{2}=6,\ x_{1}=0$ & $12$ \\
\hline
\end{tabular}
\end{center}
The point $(3;1.5)$ is computed from the two equations: $x_{1}=6-2x_{2}$ substituted into
$6x_{1}+4x_{2}=24$ gives $36-8x_{2}=24$, so $x_{2}=3/2$, $x_{1}=3$, and the profit is
$15+6=21$.

\emph{Second, the dual and strong duality.} The dual minimizes $24y_{1}+6y_{2}$ subject to
$6y_{1}+y_{2}\ge5$ and $4y_{1}+2y_{2}\ge4$. Since $x_{1},x_{2}>0$, complementary slackness
requires both dual constraints to be active; solving both as equalities gives $y_{1}=3/4$ and
$y_{2}=1/2$. The value is $18+3=21$, the same as the profit. A cross-check:
$6\cdot\tfrac34+\tfrac12=5$ and $4\cdot\tfrac34+2\cdot\tfrac12=4$, exactly the right-hand
sides, so the gap is zero.

\emph{Third, shadow prices and their range of validity.} Perturbing the right-hand side gives
the following table, with $z$ recomputed from scratch in every row.
\begin{center}
\begin{tabular}{lccc}
\hline
perturbation & optimal $\mathbf{x}$ & $z$ & change \\
\hline
$b=(25;\,6)$ & $(3.25;\,1.375)$ & $21.75$ & $+0.75$ \\
$b=(24;\,7)$ & $(2.5;\,2.25)$ & $21.50$ & $+0.50$ \\
$b=(12;\,6)$ & $(0;\,3)$ & $12.00$ & $-9.00$ \\
$b=(24;\,4)$ & $(4;\,0)$ & $20.00$ & $-1.00$ \\
\hline
\end{tabular}
\end{center}
The third row reduces the plywood to $12$, with change
$-9=0.75\cdot(-12)$; the fourth row reduces the hours by $2$, with change
$-1=0.5\cdot(-2)$. Both are still inside the intervals $b_{1}\in[12,36]$ and
$b_{2}\in[4,12]$, so the linearity holds all the way to the end. Outside them the slope changes:
with $b_{1}=40$, the relaxation gives $(0;\,3)$ with value $12$, not
$21+0.75\cdot16=33$.

\emph{Fourth, the integer program of the shop.} The enumeration of the two directions above gave
the answer $20$ at $(4;0)$, whereas its relaxation is $21$ at $(3;1.5)$. The gap is $1$, namely
$1/21=4.76\%$ of the relaxation value. Rounding down gives $19$, so it loses
$1$ unit of profit, or $5\%$ of the true answer, and that happens on a problem with only two
variables.

\emph{Fifth, a knapsack with four items.} Four items with weights $5,4,3,2$ and
values $10,7,5,3$ are loaded into a container of capacity $8$. Their value densities are $2$,
$1.75$, $1.667$, $1.5$, so the greedy rule takes the first item and then $\tfrac34$ of the second:
weight $5+3=8$, value $15.25$. The best integer answer is the first and the third item, weight
$8$ exactly full, value $15$. Its integrality gap is $0.25$, only $1.64\%$ of the bound. But
rounding down throws the second item away entirely and leaves the value $10$, namely
$33.3\%$ below the true answer: a small gap does not mean that rounding is safe, since the two
measure different things.

\kertas{For the shop above, suppose the assembly machine is available for only $5$ hours, not
$6$. Recompute the vertices, determine the optimal continuous profit, then determine the optimal
integer profit with a small enumeration. Before computing, write down your guess: does the
integrality gap grow or shrink? After that, inspect the shadow price of assembly hours at
$b_{2}=5$ and compare it with $0.5$. As an extra, find the $b_{2}$ at which the best integer
answer moves away from $(4;0)$.}

\section{Computational lab}
\tingkat{advanced}

The two blocks below are run with Python 3.14.4, numpy 2.5.2, and scipy 1.18.0. The first block
solves the linear program, its dual, and the perturbations of the right-hand side. The second
block solves the branch-and-bound tree, the knapsack, and the vertex cover of the complete graph
on five vertices. The output is copied as it stands.

\begin{lab}
\begin{verbatim}
import numpy as np
from scipy.optimize import linprog

np.set_printoptions(precision=6, suppress=True, linewidth=78)

c = np.array([5., 4.])              # profit per shelf, per bench
A = np.array([[6., 4.], [1., 2.]])  # plywood, assembly hours
b = np.array([24., 6.])

r = linprog(-c, A_ub=A, b_ub=b, bounds=[(0, None)]*2, method="highs")
print("LP   x* =", r.x, "  z* =", -r.fun)
print("residual b - Ax =", b - A @ r.x)
print("marginals (scipy):", r.ineqlin.marginals)

rd = linprog(b, A_ub=-A.T, b_ub=-c, bounds=[(0, None)]*2, method="highs")
print("dual y* =", rd.x, "  dual value =", rd.fun)
print("duality gap =", rd.fun + r.fun)

print("B^-1 =", np.linalg.inv(A), "  x_B =", np.linalg.inv(A) @ b)

for nama, db in [("b=(25, 6)", (1., 0.)), ("b=(24, 7)", (0., 1.)),
                 ("b=(12, 6)", (-12., 0.)), ("b=(24, 4)", (0., -2.))]:
    rr = linprog(-c, A_ub=A, b_ub=b + np.array(db), bounds=[(0, None)]*2,
                 method="highs")
    print(f"{nama:<10} x = {np.round(rr.x, 6) + 0.0}  z = {-rr.fun:7.4f}"
          f"  delta z = {-rr.fun + r.fun:+.4f}")
\end{verbatim}
The actual output:
\begin{verbatim}
LP   x* = [3.  1.5]   z* = 21.0
residual b - Ax = [0. 0.]
marginals (scipy): [-0.75 -0.5 ]
dual y* = [0.75 0.5 ]   dual value = 21.0
duality gap = 0.0
B^-1 = [[ 0.25  -0.5  ]
 [-0.125  0.75 ]]   x_B = [3.  1.5]
b=(25, 6)  x = [3.25  1.375]  z = 21.7500  delta z = +0.7500
b=(24, 7)  x = [2.5  2.25]  z = 21.5000  delta z = +0.5000
b=(12, 6)  x = [0. 3.]  z = 12.0000  delta z = -9.0000
b=(24, 4)  x = [4. 0.]  z = 20.0000  delta z = -1.0000
\end{verbatim}
\end{lab}

\begin{lab}
\begin{verbatim}
import itertools
import numpy as np
from scipy.optimize import linprog, milp, LinearConstraint, Bounds

c = np.array([5., 4.])
A = np.array([[6., 4.], [1., 2.]])
b = np.array([24., 6.])


def lp_bengkel(tambahan_A=None, tambahan_b=None):
    Aub = A if tambahan_A is None else np.vstack([A, tambahan_A])
    bub = b if tambahan_b is None else np.concatenate([b, tambahan_b])
    r = linprog(-c, A_ub=Aub, b_ub=bub, bounds=[(0, None)]*2, method="highs")
    if r.status != 0:
        return (None, "infeasible")
    return (-r.fun, np.round(r.x, 6) + 0.0)


print("-- branch and bound tree --")
simpul = [("root",        None, None),
          ("x2<=1",       [[0., 1.]],  [1.]),
          ("x2>=2",       [[0., -1.]], [-2.]),
          ("x2<=1,x1<=3", [[0., 1.], [1., 0.]],  [1., 3.]),
          ("x2<=1,x1>=4", [[0., 1.], [-1., 0.]], [1., -4.])]
for nama, ta, tb in simpul:
    z, x = lp_bengkel(None if ta is None else np.array(ta),
                      None if tb is None else np.array(tb))
    print(f"{nama:<14} z_LP = {z:8.4f}  x = {x}")

rb = milp(-c, constraints=LinearConstraint(A, -np.inf, b),
          integrality=[1, 1], bounds=Bounds(0, np.inf))
print("IP optimum   z_IP =", -rb.fun, " x =", np.round(rb.x, 6) + 0.0,
      " gap =", 21.0 + rb.fun)

print("-- knapsack, 4 items, capacity 8 --")
w = np.array([5., 4., 3., 2.])
v = np.array([10., 7., 5., 3.])
for mask in range(16):
    S = [i + 1 for i in range(4) if mask >> i & 1]
    bw, bv = w[[i - 1 for i in S]].sum(), v[[i - 1 for i in S]].sum()
    print(f"{str(S):<12} weight {bw:4.0f}  value {bv:4.0f}"
          f"  {'feasible' if bw <= 8 else 'infeasible'}")
rk = milp(-v, constraints=LinearConstraint(w, -np.inf, 8.),
          integrality=[1]*4, bounds=Bounds(0, 1))
print("enumeration/MILP z_IP =", -rk.fun, " x =", np.round(rk.x, 6) + 0.0)
rl = linprog(-v, A_ub=[w], b_ub=[8.], bounds=[(0, 1)]*4, method="highs")
print("relaxation      z_LP =", -rl.fun, " x =", np.round(rl.x, 6) + 0.0)
print("integrality gap =", -rl.fun + rk.fun)

print("-- vertex cover of K5: LP vs IP --")
n = 5
E = list(itertools.combinations(range(n), 2))
M = np.zeros((len(E), n))
for k, (i, j) in enumerate(E):
    M[k, i] = M[k, j] = 1.
rv = linprog(np.ones(n), A_ub=-M, b_ub=-np.ones(len(E)),
             bounds=[(0, 1)]*n, method="highs")
rc = milp(np.ones(n), constraints=LinearConstraint(M, 1, np.inf),
          integrality=[1]*n, bounds=Bounds(0, 1))
print("LP =", rv.fun, " x =", np.round(rv.x, 6) + 0.0)
print("IP =", rc.fun, " x =", np.round(rc.x, 6) + 0.0)
\end{verbatim}
The actual output:
\begin{verbatim}
-- branch and bound tree --
root           z_LP =  21.0000  x = [3.  1.5]
x2<=1          z_LP =  20.6667  x = [3.333333 1.      ]
x2>=2          z_LP =  18.0000  x = [2. 2.]
x2<=1,x1<=3    z_LP =  19.0000  x = [3. 1.]
x2<=1,x1>=4    z_LP =  20.0000  x = [4. 0.]
IP optimum   z_IP = 20.0  x = [4. 0.]  gap = 1.0
-- knapsack, 4 items, capacity 8 --
[]           weight    0  value    0  feasible
[1]          weight    5  value   10  feasible
[2]          weight    4  value    7  feasible
[1, 2]       weight    9  value   17  infeasible
[3]          weight    3  value    5  feasible
[1, 3]       weight    8  value   15  feasible
[2, 3]       weight    7  value   12  feasible
[1, 2, 3]    weight   12  value   22  infeasible
[4]          weight    2  value    3  feasible
[1, 4]       weight    7  value   13  feasible
[2, 4]       weight    6  value   10  feasible
[1, 2, 4]    weight   11  value   20  infeasible
[3, 4]       weight    5  value    8  feasible
[1, 3, 4]    weight   10  value   18  infeasible
[2, 3, 4]    weight    9  value   15  infeasible
[1, 2, 3, 4] weight   14  value   25  infeasible
enumeration/MILP z_IP = 15.0  x = [1. 0. 1. 0.]
relaxation      z_LP = 15.25  x = [1.   0.75 0.   0.  ]
integrality gap = 0.25
-- vertex cover of K5: LP vs IP --
LP = 2.5  x = [0.5 0.5 0.5 0.5 0.5]
IP = 4.0  x = [1. 1. 1. 1. 0.]
\end{verbatim}
\end{lab}

Three things must be read from the first output. The reported $\mathbf{x}^{*}$,
$(3;\,1.5)$, with $z^{*}$ exactly $21$ and no residual in either constraint, matches
complementary slackness. \texttt{scipy} reports the multipliers $(-0.75,\,-0.5)$, the negative
of the shadow prices, because the routine minimizes $-\mathbf{c}^{\top}\mathbf{x}$; the true
numbers are $0.75$ and $0.5$, and the way to be sure is to recompute through $B^{-1}$. The
duality gap comes out exactly $0$, not $10^{-15}$, because $21$ is represented exactly in
binary; on decimal data a residual of order $10^{-12}$ is normal.

The second output matches the hand computation: five nodes with the same values to four
decimals, the integer answer $20$ at $(4;\,0)$, the knapsack relaxation $15.25$, and its integer
answer $15$ with gap $0.25$. The enumeration of all sixteen subsets confirms that $15$ is the
best: every subset of higher value exceeds the capacity.

The last part of the second output checks the factor-two claim above. For the complete graph on
five vertices, the relaxation is $2.5$ whereas the integer answer is $4$; the ratio is $1.6$,
and for $n$ vertices the ratio $2(n-1)/n$ approaches $2$ from below. The relaxation bound here
can be off by nearly twice.

\begin{untukdosen}
The order that works: start from the shop and the half bench, then standard form, then vertices,
and only then the simplex. Do not start from the tableau. Students who see the tableau first
memorize the pivot step without knowing why that step may stop; those who have already seen that
the answer is at a corner will treat a pivot as a way of moving from corner to corner, and the
optimality test as a way of knowing that no corner is better.

Two typical misconceptions. First, ``the vertices are all the intersections of the constraint
lines''. Many intersections lie outside the feasible region, for example $(4;3)$ in our example,
which cuts both lines but violates both; checking the feasibility of every candidate must be
done on the board. Second, ``the integer answer is the rounding of the fractional answer''. The
shop gives $19$ against $20$, the knapsack gives $10$ against $15$; both are enough to break
that guess.

Class activity: write the initial tableau on the board, ask the class to choose the entering
column and compute the ratio test, then predict the profit before the pivot. After the final
tableau, ask for the profit when the plywood becomes $25$ sheets; $21.75$ opens the way to the
interval $[12,36]$.

Sign of readiness to go on: students can explain why the simplex method does not need to inspect
every vertex, and can name one reason why the relaxation must not be rounded.
\end{untukdosen}

\section{Summary}

\begin{ringkasan}
\begin{itemize}
  \item A linear program demands an objective and constraints of straight-line shape. The
        standard form $\max\mathbf{c}^{\top}\mathbf{x}$ with $A\mathbf{x}=\mathbf{b}$ and
        $\mathbf{x}\ge\mathbf{0}$ covers every case, with slack variables to turn inequalities
        into equalities.
  \item If the feasible region is bounded and the objective is bounded above, the best value is
        attained at a vertex. In standard form, a vertex is the same as a basic feasible
        solution, and their number is $\binom{n}{m}$, so inspecting all of them is not a
        sensible plan.
  \item The simplex method moves from vertex to vertex through columns with negative reduced
        cost and a ratio test, and stops when there is no negative reduced cost; at that moment
        the $z$ row already contains the dual prices.
  \item Weak duality holds without conditions; strong duality states that the two values are
        equal when both are finite; complementary slackness states that a loose constraint has
        price zero.
  \item The shadow price $y_{i}$ is the rate of change of the optimal value with respect to the
        right-hand side, and it holds only as long as the set of active constraints does not
        change; in the shop $0.75$ per sheet and $0.5$ per hour, valid for $b_{1}\in[12,36]$
        and $b_{2}\in[4,12]$.
  \item The LP relaxation gives a valid bound, not an answer. Rounding can give an infeasible
        point, such as $(3;2)$, or a feasible point that is not the best, such as $(3;1)$ with
        value $19$ against $20$.
  \item Branch and bound uses the relaxation bound to discard branches. The shop tree has five
        nodes, its integer answer is $20$ with gap $1$. The worst-case time is still
        exponential, and the decision version is NP-complete.
\end{itemize}
\end{ringkasan}

\section{Exercises}

\begin{latihan}
\begin{enumerate}
  \item A bakery uses $10$ kg of flour and $8$ oven hours. Cake A uses $2$ kg and
        $1$ hour with profit $30$; cake B uses $1$ kg and $2$ hours with profit $20$.
        Write the linear program, draw its feasible region, and compute the profit at every
        vertex. Then determine the best integer answer by enumeration, and compare it
        with the result of rounding down the continuous answer.

  \item Prove that the set $\{\mathbf{x}:A\mathbf{x}\le\mathbf{b}\}$ is convex, then
        explain why that proof does not apply to the set
        $\{\mathbf{x}\in\mathbb{Z}^{n}:A\mathbf{x}\le\mathbf{b}\}$.

  \item Change the profit of the shop to $x_{1}+2x_{2}$. Compute its value at the four vertices,
        state which vertex is optimal, and show that every point of the segment between the two
        optimal vertices is also optimal. Explain why a point in the middle of that segment is
        not a vertex, then connect the result with the shadow price of the hours constraint.

  \item Solve with the simplex tableau: maximize $3x_{1}+2x_{2}$ subject to
        $x_{1}+x_{2}\le4$, $x_{1}\le3$, $x_{2}\le3$, $x\ge0$. Write every intermediate tableau,
        name the entering column and the reason for it, then read the dual prices from the final
        tableau. (Hint: two pivots suffice; compare the final result with the vertex obtained
        from the intersection of $x_{1}+x_{2}=4$ and $x_{1}=3$.)

  \item Write the dual of the program in exercise 1, solve it, and check strong duality
        numerically. Explain the meaning of every shadow price in one sentence.

  \item Assess this claim: \emph{rounding down the answer of the relaxation always gives a
        feasible integer point.} Decide true, false, or not enough evidence, and
        exhibit one counterexample with numbers.

  \item Assess this claim: \emph{a positive shadow price means that the constraint must be
        binding.} Decide true, false, or not enough evidence, then explain the actual relation
        between shadow prices and complementary slackness.
        (Hint: inspect what happens to a constraint with $y_{i}>0$ according to
        $y_{i}(b_{i}-\mathbf{a}_{i}^{\top}\mathbf{x})=0$.)

  \item For the knapsack with weights $(5,4,3,2)$, values $(10,7,5,3)$, and capacity
        $8$, compute its LP relaxation by hand using the density order, then compare it with the
        enumeration result. Compute the integrality gap and state it as a percentage of the
        bound. (Hint: sort by value per unit of weight, then fill the container to the top; the
        last item enters partly.)

  \item Repeat exercise 8 with capacity $7$. Determine the relaxation, the integer answer, and
        the gap, then explain why the gap can change although only one data value changed.

  \item For the vertex cover of the complete graph on $n$ vertices, prove that the relaxation has
        value $n/2$ and the integer answer $n-1$, then compute the ratio of the two for
        $n=4,5,10$. Explain why that result shows that rounding is not a general way to solve
        integer programs.
\end{enumerate}
\end{latihan}

\begin{untukpeneliti}
What is really proved in this chapter: weak duality for a primal-dual pair in standard form,
without any condition; the converse direction of complementary slackness as an optimality test;
optimality at a vertex for a bounded feasible region, with a proof through the active-constraint
argument; and the validity of pruning in branch and bound. What is borrowed: strong duality,
whose usual proof goes through Farkas' lemma or through the termination of the simplex; both are
recorded in \texttt{WSPrinted.pdf}, in the section on linear programming duality, and in Boyd and
Vandenberghe, section 5.2.

What is popularized beyond its proof is quite a lot, and three things must be named.
``The LP relaxation is nearly right'': its worst-case approximation factor can grow without
bound on certain classes of problems, and even for the vertex cover it reaches $2$. ``The shadow
price is the derivative of the optimal value'': that is true only when the optimal value is
differentiable with respect to the right-hand side; at a degenerate point, where two optimal
bases share the same vertex, the derivative does not exist and what remains is only a
subgradient. ``Branch and bound is always fast'': its worst-case time is exponential, and
slowness on a particular instance is not a sign of a bad implementation.

Open questions: how strong the Chvatal-Gomory cutting-plane hierarchy is for particular classes
of polytopes, and whether a superpolynomial lower bound exists for the size of extended
formulations of the traveling salesman polytope; the results of Fiorini and coauthors in 2012
show an exponential bound for that case, while the search for similar bounds for other problems
continues. Limits of validity of this chapter: rational data, exact arithmetic for proofs,
floating-point tolerance in practice; all statements about shadow prices hold for the continuous
model, not for the integer model.

\peta{This chapter follows \emph{Introduction to Linear Optimization} in the file
\texttt{WSPrinted.pdf}: the sections on formulation, the geometry of polyhedral sets (vertices
and extreme points), systems of linear inequalities and duality (complementary slackness,
sensitivity analysis), and the simplex method (basic feasible solutions, reduced costs, the
optimality test). The integer-programming side follows Wolsey, \emph{Integer Programming},
chapter 1 (formulation) and chapter 2 (optimality, relaxation, bounds), which places the
relaxation as the source of bounds and branching as the way to use them. The convex viewpoint in
the second and fourth parts is Boyd and Vandenberghe, \emph{Convex Optimization}, chapters 4 and
5. Relation to other chapters: standard form and the feasible region continue Chapter 21,
solving linear systems uses Chapter 8, and the integrality gap is used again in Chapter 25.}

\end{untukpeneliti}
